\section{Introduction}
\label{sec:introduction}

% Source-note convention: Verbatim means unchanged prose apart from formatting;
% Edited includes condensed or combined source text; New identifies report synthesis.
% Source: ../egg-papers/2026-OOPSLA-eggcc/01-introduction.tex, opening.
\sourcenote{Edited from \emph{Efficient Extraction for Effectful E-graphs}, \S1.}
E-graphs power state-of-the-art program optimizers, synthesizers, and verifiers
  across a variety of domains~\cite{egg,szalinski,caviar}.
Through \emph{equality saturation}, an optimizer applies semantics-preserving rules
  to build a space of equivalent programs, then \emph{extracts} one using a cost function~\cite{peggy}.
By retaining alternatives and separating exploration from extraction, equality saturation
  lets engineers focus on optimization rules instead of transformation order.

% Source: ../egg-papers/2020-big-paper/01-intro.tex ("Ideally...");
% remaining problem framing from the report abstract and Oliver's proposal.
\sourcenote{Edited from \emph{egg}, \S1; new problem framing from the research plan.}
Ideally, a user could simply provide a language grammar and rewrites,
  and equality saturation would produce an effective optimizer.
Realizing this ideal requires more than applying rewrites.
Applications need to justify their results, handle effects and variable binding,
  and combine rewriting with program analyses.
These requirements affect the evidence retained during search,
  the programs extraction can select, and the extensions that must work together.
A verified optimizer for effectful programs must address them together.

\sourcenote{Verbatim from the report abstract.}
This thesis proposal presents semantic extensions to e-graphs that address these challenges.
Its central claim is that \textbf{extended e-graphs are practical and beneficial
  for program reasoning, including verified compilation
  of effectful programs}.

% Sources: Small Proofs, 04-flattened-optimal.tex and 05-flattened-greedy.tex;
% eggcc, 01-introduction.tex. Shared theme; each parameter is introduced with its project below.
\sourcenote{New synthesis of the parameterized-algorithm contributions.}
These extensions must also be efficient enough to use in optimizers.
Both proof minimization and effect-safe extraction are NP-hard~\cite{eggproof,eggcc}.
We develop parameterized algorithms whose cost depends on structural properties of the input
  as well as its overall size.
Identifying and controlling these parameters guides the implementation of practical
  proof generation and extraction.

\subsection{Proof Production and Minimization}
\label{sec:intro-proofs}

% Source: ../egg-papers/2022-explanations-fmcad/01-intro.tex;
% Herbie proof documentation; condensed contribution overview.
\sourcenote{Edited from \emph{Small Proofs from Congruence Closure}, \S1;
  explanation and debugging motivation added.}
A \emph{proof certificate} records equalities that establish why two terms are equivalent.
Clients can check the certificate independently of the optimizer.
Proofs also explain transformations and help locate bugs, as in Herbie~\cite{herbie-proofs}.
Smaller proofs give users fewer steps to follow and developers fewer rule applications to inspect.
Large proofs can cost more to check than the optimization itself.

% Sources: ../egg-papers/2022-explanations-fmcad/01-intro.tex,
% 04-flattened-optimal.tex, and 05-flattened-greedy.tex.
\sourcenote{Edited from \emph{Small Proofs from Congruence Closure}, \S\S1, 4--5;
  the shared parameter emphasized.}
Our work presents two algorithms for producing smaller proofs~\cite{eggproof}:
  \textsc{TreeOpt}, a theoretical algorithm that optimizes a relaxed proof-size objective,
  and greedy proof search, the practical algorithm implemented in \texttt{egg}.
Both algorithms are parameterized by the number of \emph{congruence} justifications,
  which use equalities between subterms to justify equality between enclosing terms.
Bounding this count and the number of subproof searches makes greedy proof search practical.
The greedy algorithm reduces validation costs while retaining alternative explanations
  that traditional proof generation discards.

\subsection{Effect-Safe Extraction}
\label{sec:intro-extraction}

% Source: ../egg-papers/2026-OOPSLA-eggcc/abstract.tex and 01-introduction.tex.
\sourcenote{Edited from \emph{Efficient Extraction for Effectful E-graphs}, abstract and \S1;
  new connection from proof production to extraction.}
An optimizer must also turn the equalities it discovers into an executable program.
For programs with side effects, extraction must select a valid ordering of memory
  and input/output operations.
We prove that finding any \emph{effect-safe extraction} is NP-complete,
  then develop statewalk DP, an algorithm that enforces effect ordering without external solvers~\cite{eggcc}.
Its running time depends on \emph{statewalk width},
  which measures the dataflow interactions among effects and is usually small in our benchmarks.
The implementation in \texttt{eggcc} makes extraction practical
  and enables useful optimizations of imperative programs.

\subsection{Composable E-Graph Extensions}
\label{sec:intro-extensions}

% Source: ../extended-egglog-paper/intro.tex; submission status from Oliver.
\sourcenote{Edited from \emph{Extended E-Graphs via Compilation}, \S1;
  submission status added.}
Real optimizers combine these capabilities with analyses and variable binding.
When a rewrite depends on an analysis result, its proof must justify that result;
  when terms contain bound variables, matching and proof checking must respect renaming.
Our work currently under submission implements extensions as compiler passes
  over \texttt{egglog}, a language for programming with e-graphs~\cite{egglog,extended-egglog}.
One pass replaces built-in equality maintenance with explicit rules;
  a later proof pass records the rule applications that establish equalities and analysis facts.
This architecture combines proof production, analyses, and slotted e-graphs for bound variables.

\subsection{Integration with Verified Compilers}
\label{sec:intro-verified-compilation}

% Sources: Oliver's research plan and status update; Lean egg and Coquetier papers.
\sourcenote{New: integration challenge and proposed library from Oliver's research plan.}
These extensions supply candidate programs and certificates.
Integrating them into a compiler verified in Lean or Rocq requires connecting
  each certificate to a semantic claim about the programs it relates.
Existing integrations reconstruct \emph{proof terms}, expressions in the assistant's logic
  that encode proofs, from \texttt{egg}'s explanations~\cite{lean-egg,coquetier}.
They run equality saturation while the proof assistant processes a proof script.
The challenge is to validate new programs within a deployed compiler,
  after the compiler itself has been verified and built.

\sourcenote{New: proposed library and completion status from Oliver's research plan.}
Our proposed library exports proved theorems as \texttt{egglog} rules
  and checks the resulting certificates with a verified executable checker at compiler runtime.
This lets a verified compiler use \texttt{egglog}'s substantial unverified implementation,
  without invoking the proof assistant on each input program.
The Lean checker experiment and its soundness proof for a restricted fragment are complete.
The remaining work connects theorem export, proofs of analysis facts,
  effect-safe extraction, and runtime validation into a reusable library,
  then evaluates a compiler case study.

\sourcenote{New roadmap for this report.}
Section~\ref{sec:background} introduces the foundations.
Sections~\ref{sec:proofs}--\ref{sec:extensions} present the existing contributions;
  Section~\ref{sec:proposed-research} develops the proposed integration and evaluation,
  followed by the dissertation plan in Section~\ref{sec:timeline}.
